% \section{Sistemas acoplados espaçonave--manipulador} % modelagem_controle_espaconave_manipulador.tex
% \label{sec_modelagem_controle_espaconave_manipulador}

% Nesta seção é apresentada uma revisão da literatura relacionada à modelagem e ao controle de sistemas acoplados formados por uma espaçonave e um manipulador robótico. São destacadas as referências que tratam da modelagem cinemática e dinâmica desses sistemas, do acoplamento entre a base e o manipulador, das formulações empregadas para a obtenção das equações de movimento e das estratégias de controle aplicadas a manipuladores espaciais.

% %%%%%%%%%%%% modelagem cinematica

% A modelagem cinemática de manipuladores espaciais constitui uma etapa central para descrever o movimento relativo entre os elos, determinar a posição e a orientação do efetuador final e estabelecer as relações entre o espaço das juntas e o espaço cartesiano \cite{fonseca2019kinematics}. No contexto espacial, essa modelagem adquire complexidade adicional, pois o manipulador está acoplado a uma plataforma que, em geral, não pode ser tratada como uma base fixa convencional. Em vez disso, a formulação deve considerar a interação entre o manipulador e a espaçonave, bem como o uso consistente dos referenciais inercial, orbital e fixo ao corpo \cite{fonseca2019kinematics,fonseca2016dynamics}.

% Nesse cenário, abordagens clássicas de cinemática direta e inversa continuam desempenhando papel relevante. A convenção de Denavit--Hartenberg permanece amplamente empregada para a definição sistemática dos referenciais e para a obtenção das transformações homogêneas entre os elos, permitindo calcular a pose do efetuador final a partir das variáveis articulares \cite{fonseca2019kinematics}. Quando a cinemática inversa não admite solução analítica em forma fechada, o problema pode ser formulado como uma otimização, na qual se busca a configuração articular que minimiza o erro entre a pose desejada e a pose obtida pelo manipulador \cite{santos2022uncertainty,santos2023physics}.

% Entretanto, em manipuladores do tipo espaçonave, a formulação cinemática não se limita à descrição geométrica da cadeia serial. Como a movimentação das juntas pode induzir deslocamentos e rotações da plataforma, torna-se necessário compatibilizar a cinemática do manipulador com os movimentos da espaçonave e com os referenciais relevantes ao problema orbital \cite{fonseca2019kinematics,fonseca2016dynamics}. Essa característica é especialmente importante em operações de proximidade, captura e serviço em órbita, nas quais erros de posicionamento e orientação podem comprometer o desempenho da missão \cite{arantes2010far,fonseca2016dynamics}.

% Além das formulações analíticas tradicionais, estudos mais recentes têm explorado estratégias baseadas em aprendizado de máquina para apoiar a resolução da cinemática inversa e o planejamento de movimento. Redes neurais e métodos como máquinas de vetores de suporte vêm sendo empregados para estimar configurações articulares, quantificar incertezas e melhorar a convergência de procedimentos de otimização, inclusive em situações nas quais os efeitos dinâmicos da plataforma e do manipulador são incorporados ao problema \cite{santos2022uncertainty,santos2023physics,santos2022machine}.

% %%%%%%%%%%%%%%%%%%

% Diferentemente dos manipuladores fixos em solo, os manipuladores espaciais operam em sistemas nos quais o movimento das juntas produz reações sobre a base, alterando sua atitude e, em determinadas configurações, também seu movimento translacional.
% Dessa forma, a modelagem e o controle desses sistemas exigem o tratamento explícito do acoplamento dinâmico entre a espaçonave e o manipulador, não sendo suficiente considerar a base como um suporte rigidamente fixo. 
% Conforme estabelecido classicamente por \cite{umetani1989resolved}, \cite{papadopoulos1991nature} e \cite{dubowsky1993kinematics}, a formulação da cinemática e da dinâmica, como o uso da Jacobiana generalizada, deve incorporar a natureza flutuante da base e suas restrições de controle de forma integrada.
% No que se refere aos avanços posteriores na formulação dinâmica desses sistemas, \cite{from2011singularityfree} apresentam equações dinâmicas livres de singularidades para sistemas espaçonave--manipulador, utilizando uma representação mínima baseada em quase-coordenadas em um arcabouço lagrangiano.
% Essa abordagem permite evitar as singularidades inerentes às parametrizações de atitude por ângulos de Euler, resultando em um modelo computacionalmente mais robusto para simulações de malha fechada.

% Os trabalhos de \cite{floresabad2014review} apresentam uma revisão sobre tecnologias de robótica espacial para \textit{on-orbit servicing}, discutindo problemas de cinemática, dinâmica, controle e verificação experimental de sistemas robóticos acoplados. Os autores destacam que a modelagem e o controle de manipuladores espaciais continuam entre os principais desafios das missões autônomas em órbita, especialmente em aplicações de captura, reparo, montagem, reabastecimento e remoção de detritos. Essa revisão é importante por situar os desenvolvimentos teóricos em um contexto operacional mais amplo.

% Os estudos feitos por \cite{fonseca2017attitude} analisam a dinâmica de atitude de uma espaçonave tratada de forma análoga a um manipulador robótico no contexto de missões de \textit{on-orbit servicing}. Os autores formulam o problema com base na dinâmica rotacional em órbita e avaliam o uso de controladores LQR e PID para as formas linearizada e não linear do modelo. As simulações realizadas em MATLAB mostram que ambas as abordagens atendem aos objetivos definidos para operações de \textit{rendezvous} e \textit{berthing}, permitindo ainda comparar esforços de controle e comportamento dinâmico sob diferentes formulações. Essa referência é particularmente relevante para a presente dissertação por enfatizar a necessidade de tratar de forma integrada a atitude da espaçonave e a dinâmica associada ao sistema robótico embarcado.

% Mais recentemente, \cite{rybus2017nmpc} apresentam um sistema de controle para manipuladores espaciais \textit{free-floating} baseado em controle preditivo não linear. A arquitetura proposta combina planejamento de trajetória com um controlador NMPC, ambos formulados de modo a considerar explicitamente a natureza flutuante do sistema espaçonave--manipulador. O método é avaliado por simulações numéricas em um caso planar simplificado, incluindo perturbações e incertezas paramétricas. Os resultados apresentados indicam desempenho superior ao de estratégias baseadas na inversa da Jacobiana dinâmica e no \textit{Modified Simple Adaptive Control}. Essa referência mostra uma tendência mais recente da literatura em direção a estratégias de controle mais integradas, capazes de lidar com rastreamento, restrições dinâmicas e robustez a incertezas.

% De modo geral, a literatura revisada mostra que a modelagem e o controle de sistemas acoplados espaçonave--manipulador evoluíram de formulações cinemáticas fundamentais para representações dinâmicas mais robustas e, posteriormente, para arquiteturas de controle mais sofisticadas. Os trabalhos revisados mostram que o tratamento do acoplamento entre a espaçonave e o manipulador constitui aspecto central do problema, sendo indispensável para a descrição adequada do comportamento do sistema em operações orbitais. Esses aspectos são relevantes para a presente dissertação, pois reforçam a necessidade de formular de modo coerente a modelagem e o controle do sistema acoplado nas fases de aproximação e interação com o alvo.

\section{Sistemas acoplados espaçonave--manipulador} % modelagem_controle_espaconave_manipulador.tex
\label{sec_modelagem_controle_espaconave_manipulador}

Nesta seção é apresentada uma revisão da literatura relacionada à modelagem e ao controle de sistemas acoplados formados por uma espaçonave e um manipulador robótico. São destacadas as referências que tratam da modelagem cinemática e dinâmica desses sistemas, do acoplamento entre a base e o manipulador, das formulações empregadas para a obtenção das equações de movimento e das estratégias de controle aplicadas a manipuladores espaciais.

A modelagem cinemática de manipuladores espaciais constitui uma etapa central para descrever o movimento relativo entre os elos, determinar a posição e a orientação do efetuador final e estabelecer as relações entre o espaço das juntas e o espaço cartesiano \cite{fonseca2019kinematics}. No contexto espacial, essa modelagem adquire complexidade adicional, pois o manipulador está acoplado a uma plataforma que, em geral, não pode ser tratada como uma base fixa convencional. Em vez disso, a formulação deve considerar a interação entre o manipulador e a espaçonave, bem como o uso consistente dos referenciais inercial, orbital e fixo ao corpo \cite{fonseca2019kinematics,fonseca2016dynamics}.

Nesse cenário, abordagens clássicas de cinemática direta e inversa continuam desempenhando papel relevante. A convenção de Denavit--Hartenberg permanece amplamente empregada para a definição sistemática dos referenciais e para a obtenção das transformações homogêneas entre os elos, permitindo calcular a pose do efetuador final a partir das variáveis articulares \cite{fonseca2019kinematics}. Quando a cinemática inversa não admite solução analítica em forma fechada, o problema pode ser formulado como uma otimização, na qual se busca a configuração articular que minimiza o erro entre a pose desejada e a pose obtida pelo manipulador \cite{santos2022uncertainty,santos2023physics}.

Entretanto, em manipuladores do tipo espaçonave, a formulação cinemática não se limita à descrição geométrica da cadeia serial. Como a movimentação das juntas pode induzir deslocamentos e rotações da plataforma, torna-se necessário compatibilizar a cinemática do manipulador com os movimentos da espaçonave e com os referenciais relevantes ao problema orbital \cite{fonseca2019kinematics,fonseca2016dynamics}. Essa característica é especialmente importante em operações de proximidade, captura e serviço em órbita, nas quais erros de posicionamento e orientação podem comprometer o desempenho da missão \cite{arantes2010far,fonseca2016dynamics}.

Além das formulações analíticas tradicionais, estudos mais recentes têm explorado estratégias baseadas em aprendizado de máquina para apoiar a resolução da cinemática inversa e o planejamento de movimento. Redes neurais e métodos como máquinas de vetores de suporte vêm sendo empregados para estimar configurações articulares, quantificar incertezas e melhorar a convergência de procedimentos de otimização, inclusive em situações nas quais os efeitos dinâmicos da plataforma e do manipulador são incorporados ao problema \cite{santos2022uncertainty,santos2023physics,santos2022machine}.

Do ponto de vista dinâmico, os manipuladores espaciais diferem dos manipuladores fixos em solo, pois operam em sistemas nos quais o movimento das juntas produz reações sobre a base, alterando sua atitude e, em determinadas configurações, também seu movimento translacional. Dessa forma, a modelagem e o controle desses sistemas exigem o tratamento explícito do acoplamento entre a espaçonave e o manipulador, não sendo suficiente considerar a base como um suporte rigidamente fixo. Conforme estabelecido classicamente por \cite{umetani1989resolved}, \cite{papadopoulos1991nature} e \cite{dubowsky1993kinematics}, a formulação da cinemática e da dinâmica, incluindo o uso da Jacobiana generalizada, deve incorporar, de forma integrada, a natureza flutuante da base e as restrições de controle associadas a esse tipo de sistema.

No que se refere aos avanços posteriores na formulação dinâmica desses sistemas, \cite{from2011singularityfree} apresentam equações dinâmicas livres de singularidades para sistemas espaçonave--manipulador, utilizando uma representação mínima baseada em quase-coordenadas em um arcabouço lagrangiano. Essa abordagem permite evitar as singularidades inerentes às parametrizações de atitude por ângulos de Euler, resultando em um modelo computacionalmente mais robusto para simulações em malha fechada.

Os trabalhos de \cite{floresabad2014review} apresentam uma revisão sobre tecnologias de robótica espacial para \textit{on-orbit servicing}, discutindo problemas de cinemática, dinâmica, controle e verificação experimental de sistemas robóticos acoplados. Os autores destacam que a modelagem e o controle de manipuladores espaciais continuam entre os principais desafios das missões autônomas em órbita, especialmente em aplicações de captura, reparo, montagem, reabastecimento e remoção de detritos. Essa revisão é importante por situar os desenvolvimentos teóricos em um contexto operacional mais amplo.

Os estudos de \cite{fonseca2017attitude} analisam a dinâmica de atitude de uma espaçonave tratada de forma análoga a um manipulador robótico no contexto de missões de \textit{on-orbit servicing}. Os autores formulam o problema com base na dinâmica rotacional em órbita e avaliam o uso de controladores \gls{lqr} e \gls{pid} para as formas linearizada e não linear do modelo. As simulações realizadas em MATLAB mostram que ambas as abordagens atendem aos objetivos definidos para operações de \textit{rendezvous} e \textit{berthing}, permitindo ainda comparar esforços de controle e comportamento dinâmico sob diferentes formulações. Essa referência é particularmente relevante para a presente dissertação, pois enfatiza a necessidade de tratar de forma integrada a atitude da espaçonave e a dinâmica associada ao sistema robótico embarcado.

Mais recentemente, \cite{rybus2017nmpc} apresentam um sistema de controle para manipuladores espaciais \textit{free-floating} baseado em controle preditivo não linear. A arquitetura proposta combina planejamento de trajetória com um controlador \gls{nmpc}, ambos formulados de modo a considerar explicitamente a natureza flutuante do sistema espaçonave--manipulador. O método é avaliado por simulações numéricas em um caso planar simplificado, incluindo perturbações e incertezas paramétricas. Os resultados apresentados indicam desempenho superior ao de estratégias baseadas na inversa da Jacobiana dinâmica e no \textit{Modified Simple Adaptive Control}. Essa referência mostra uma tendência mais recente da literatura em direção a estratégias de controle mais integradas, capazes de lidar com rastreamento, restrições dinâmicas e robustez a incertezas.

De modo geral, a literatura revisada mostra que a modelagem e o controle de sistemas acoplados espaçonave--manipulador evoluíram de formulações cinemáticas fundamentais para representações dinâmicas mais robustas e, posteriormente, para arquiteturas de controle mais sofisticadas. Os trabalhos revisados mostram que o tratamento do acoplamento entre a espaçonave e o manipulador constitui um aspecto central do problema, sendo indispensável para a descrição adequada do comportamento do sistema em operações orbitais. Esses aspectos são relevantes para a presente dissertação, pois reforçam a necessidade de formular, de modo coerente, a modelagem e o controle do sistema acoplado nas fases de aproximação e interação com o alvo.